\documentclass[12pt]{article}
\usepackage[utf8]{inputenc}
\usepackage[T1]{fontenc}
\usepackage[margin=1in]{geometry}
\usepackage{setspace}
\usepackage{xcolor}
\usepackage[colorlinks=true, urlcolor=blue!60!black, linkcolor=blue!60!black]{hyperref}
\setlength{\parindent}{2em}
\setlength{\parskip}{1.25em}
\renewcommand{\baselinestretch}{1.0}
\title{\textit{United States v. Vaello Madero}, 596 U.S. 159 (2022)}
\author{Case Brief}
\date{}
\begin{document}
\maketitle

\section{Case Metadata}
\begin{itemize}
  \item \textbf{Case Name:} \emph{United States v. Vaello Madero}, \href{https://supreme.justia.com/cases/federal/us/596/20-303/}{596 U.S. 159 (2022)}
  \item \textbf{Docket No.:} 20-303 \quad \textbf{Argued:} November 9, 2021 \quad \textbf{Decided:} April 21, 2022
  \item \textbf{Court:} Supreme Court of the United States (on writ of certiorari to the U.S. Court of Appeals for the First Circuit, 956 F.3d 12 (1st Cir. 2020))
  \item \textbf{Briefing Book Placement:} Chapter 3: Federal Legislative Power
  \item \textbf{Vote:} 8--1 (reversing the First Circuit)
  \item \textbf{Opinion of the Court:} Justice Kavanaugh (joined by Chief Justice Roberts and Justices Thomas, Breyer, Alito, Kagan, Gorsuch, and Barrett)
  \item \textbf{Concurring Opinions:} Justice Thomas; Justice Gorsuch
  \item \textbf{Dissenting Opinion:} Justice Sotomayor
  \item \textbf{Key Constitutional and Statutory Provisions:} Territory Clause, U.S. Const. art. IV, \S~3, cl. 2; Fifth Amendment Due Process Clause (equal-protection component); Fourteenth Amendment Citizenship Clause (discussed in concurrence); Social Security Act Supplemental Security Income (SSI) provisions, 42 U.S.C. \S\S~1382c(a)(1)(B)(i), (e); Puerto Rico federal tax exemptions, 48 U.S.C. \S~734, 26 U.S.C. \S~933
  \item \textbf{Primary Source Links:} \href{https://www.law.cornell.edu/supremecourt/text/20-303}{Cornell LII Supreme Court Collection}; \href{https://supreme.justia.com/cases/federal/us/596/20-303/}{Justia U.S. Supreme Court Center}
\end{itemize}

\section{Facts and Procedural History}
Jose Luis Vaello Madero, a U.S. citizen born in Puerto Rico in 1954, moved to New York in 1985. In 2012, after suffering a severe illness while residing in New York, he qualified for and began receiving monthly Supplemental Security Income (SSI) disability payments---a means-tested federal safety-net program enacted in 1972 for low-income aged, blind, and disabled individuals (\href{https://www.law.cornell.edu/supremecourt/text/20-303}{Cornell LII}). By statute, SSI eligibility is restricted to a ``resident of the United States,'' defined as the 50 States and the District of Columbia (42 U.S.C. \S\S~1382c(a)(1)(B)(i), (e)), and later extended by covenant to the Northern Mariana Islands, leaving Puerto Rico residents under a less generous, capped federal-territorial matching program known as Aid to the Aged, Blind, and Disabled (AABD).

In 2013, Vaello Madero moved back to Puerto Rico. Unaware of his change of residence, the Social Security Administration (SSA) continued depositing monthly SSI payments of \$733 to \$808 into his bank account. In June 2016, when Vaello Madero visited an SSA office in Puerto Rico to apply for Title II Social Security benefits, the agency learned of his relocation, terminated his SSI benefits retroactively to August 2013 on the ground that he was ``outside of the United States,'' and in 2017 sued him in federal court to recover \$28,081 in overpayments (\href{https://www.law.cornell.edu/supremecourt/text/20-303}{Cornell LII}).

Vaello Madero raised an affirmative constitutional defense, arguing that Congress's exclusion of Puerto Rico residents from SSI violated the equal-protection component of the Fifth Amendment's Due Process Clause. The U.S. District Court for the District of Puerto Rico granted summary judgment to Vaello Madero (356 F. Supp. 3d 208 (D.P.R. 2019)), and the First Circuit unanimously affirmed (956 F.3d 12 (1st Cir. 2020)), concluding that excluding needy citizens in Puerto Rico failed rational-basis review because SSI recipients are by definition too poor to pay federal income taxes. The Supreme Court granted certiorari (592 U.S. 1169 (2021)) and reversed.

\section{Issue(s) Presented}
Does the equal-protection component of the Fifth Amendment's Due Process Clause require Congress, when exercising its legislative authority under the Territory Clause (Article IV, \S~3, cl. 2), to make Supplemental Security Income (SSI) benefits available to residents of Puerto Rico to the same extent that Congress makes them available to residents of the 50 States?

\section{Holding}
\textbf{No.} In an 8--1 decision authored by Justice Kavanaugh, the Supreme Court held that the Constitution does not require Congress to extend SSI benefits to residents of Puerto Rico. Under deferential rational-basis review, Puerto Rico's general exemption from most federal income, gift, estate, and excise taxes supplies a rational basis for Congress to distinguish Puerto Rico residents from residents of the States in structuring federal benefits programs (\href{https://www.law.cornell.edu/supremecourt/text/20-303}{Cornell LII}).

\section{Reasoning}
Justice Kavanaugh's concise majority opinion rested on constitutional text, historical practice, controlling precedent, and structural pragmatism.

\textbf{1. Broad Authority Under the Territory Clause.} The Territory Clause empowers Congress to ``make all needful Rules and Regulations respecting the Territory . . . belonging to the United States'' (Art. IV, \S~3, cl. 2). Since acquiring Puerto Rico in 1898, Congress has exercised that broad discretion to balance national objectives with territorial autonomy, maintaining distinct tax and benefits regimes. While Puerto Rico residents pay Social Security, Medicare, and federal unemployment taxes (and receive those corresponding benefits), they are generally exempt from federal income, gift, estate, and excise taxes on Puerto Rico-source income (48 U.S.C. \S~734; 26 U.S.C. \S~933) (\href{https://www.law.cornell.edu/supremecourt/text/20-303}{Cornell LII}).

\textbf{2. Precedent Dictating Deferential Rational-Basis Review.} The Court held that two per curiam precedents controlled: \href{https://supreme.justia.com/cases/federal/us/435/1/}{\emph{Califano v. Torres}, 435 U.S. 1 (1978)} (upholding Puerto Rico's exclusion from SSI against a right-to-travel challenge), and \href{https://supreme.justia.com/cases/federal/us/446/651/}{\emph{Harris v. Rosario}, 446 U.S. 651 (1980)} (upholding lower AFDC block-grant reimbursements to Puerto Rico under Fifth Amendment equal protection). Applying deferential rational-basis review, the Court concluded that Congress may reasonably consider the ``general balance of benefits to and burdens on the residents of Puerto Rico'' without performing a ``dollar-to-dollar comparison'' at either the individual or collective level (\href{https://www.law.cornell.edu/supremecourt/text/20-303}{Cornell LII}).

\textbf{3. Structural Consequences and Legislative Discretion.} Kavanaugh warned that mandating constitutional symmetry on the benefits side would invite demands for identical federal taxation on the tax side, potentially imposing severe new financial burdens on Puerto Rico's economy. Emphasizing institutional restraint, the Court noted that Congress remains free to extend SSI to Puerto Rico by statute---as the President supported as a matter of policy---and added in Footnote 1 that its ruling did not imply Congress could exclude residents of individual States from federal benefits programs (\href{https://www.law.cornell.edu/supremecourt/text/20-303}{Cornell LII}).

\section{Concurring and Dissenting Opinions}
\begin{itemize}
  \item \textbf{Justice Thomas, concurring:} Joined the majority in full, but wrote separately to question the originalist pedigree of ``reverse incorporation'' under \href{https://supreme.justia.com/cases/federal/us/347/497/}{\emph{Bolling v. Sharpe}, 347 U.S. 497 (1954)} (and \href{https://supreme.justia.com/cases/federal/us/515/200/}{\emph{Adarand Constructors, Inc. v. Pe\~na}, 515 U.S. 200 (1995)}), which located an equal-protection guarantee inside the Fifth Amendment's Due Process Clause. Reviewing Reconstruction-era debates and Justice Harlan's \emph{Plessy} dissent, Thomas proposed that the Fourteenth Amendment's \textbf{Citizenship Clause}---not substantive due process---may supply the firmer textual foundation prohibiting federal racial discrimination in civil rights (\href{https://www.law.cornell.edu/supremecourt/text/20-303}{Cornell LII}).
  \item \textbf{Justice Gorsuch, concurring:} Joined the majority because neither party asked to overrule the \emph{Insular Cases} (assuming equal protection applied as a ``fundamental'' right in Puerto Rico), but delivered a blistering indictment of the \emph{Insular Cases} (\emph{Downes v. Bidwell}, \emph{Balzac v. Porto Rico}) as a colonial invention grounded in ``ugly racial stereotypes'' and social Darwinism. Mocking the absurdity that full constitutional protections apply on uninhabited Palmyra Atoll (``incorporated'') while 3 million U.S. citizens in Puerto Rico lack constitutional jury-trial rights (``unincorporated''), Gorsuch urged the Court to squarely overrule the \emph{Insular Cases} at the next opportunity (\href{https://www.law.cornell.edu/supremecourt/text/20-303}{Cornell LII}).
  \item \textbf{Justice Sotomayor, dissenting:} Argued that excluding needy U.S. citizens in Puerto Rico from SSI fails even rational-basis review. She distinguished \emph{Torres} and \emph{Rosario} as summary dispositions involving different claims and resting on the false premise that Puerto Rico residents pay no federal taxes (in reality contributing \$3--4 billion annually in payroll and other federal taxes). Because SSI is a uniform, nationally administered, direct-to-individual program of last resort for people too poor to owe federal income taxes anywhere, tying eligibility to territorial tax exemptions is ``antithetical to the entire premise of the program.'' She warned that the Court's logic would let Congress exclude residents of low-tax-revenue States like Vermont or Wyoming, noted that Puerto Ricans lack voting representation in Congress to fix the disparity, and agreed in Footnote 4 with Justice Gorsuch that the \emph{Insular Cases} should be overruled (\href{https://www.law.cornell.edu/supremecourt/text/20-303}{Cornell LII}).
\end{itemize}

\section{Key Quotations}
\begin{quote}
  ``In devising tax and benefits programs, it is reasonable for Congress to take account of the general balance of benefits to and burdens on the residents of Puerto Rico. In doing so, Congress need not conduct a dollar-to-dollar comparison of how its tax and benefits programs apply in the States as compared to the Territories, either at the individual or collective level.''
  --- Justice Kavanaugh, \href{https://www.law.cornell.edu/supremecourt/text/20-303}{596 U.S. at 165--66 (Cornell LII)}.
\end{quote}
\begin{quote}
  ``Rather than continue to invoke the Fifth Amendment's Due Process Clause to justify \emph{Bolling}, in an appropriate case, we should more carefully consider whether this interpretation of the Citizenship Clause would yield a similar, and more supportable, result.''
  --- Justice Thomas, concurring, \href{https://www.law.cornell.edu/supremecourt/text/20-303}{596 U.S. at 179 (Cornell LII)}.
\end{quote}
\begin{quote}
  ``It is past time to acknowledge the gravity of this error and admit what we know to be true: The \emph{Insular Cases} have no foundation in the Constitution and rest instead on racial stereotypes. They deserve no place in our law.''
  --- Justice Gorsuch, concurring, \href{https://www.law.cornell.edu/supremecourt/text/20-303}{596 U.S. at 180 (Cornell LII)}.
\end{quote}
\begin{quote}
  ``In my view, there is no rational basis for Congress to treat needy citizens living anywhere in the United States so differently from others. To hold otherwise, as the Court does, is irrational and antithetical to the very nature of the SSI program and the equal protection of citizens guaranteed by the Constitution.''
  --- Justice Sotomayor, dissenting, \href{https://www.law.cornell.edu/supremecourt/text/20-303}{596 U.S. at 190 (Cornell LII)}.
\end{quote}

\section{Doctrinal Significance and Modern Legacy (Through 2026)}
\emph{Vaello Madero} reaffirmed Congress's sweeping plenary discretion under the Territory Clause to craft asymmetric fiscal and social-welfare arrangements for U.S. Territories, locking in a highly deferential rational-basis baseline whenever geographic territorial classifications are challenged under Fifth Amendment equal protection. Following the ruling, the Congressional Budget Office explicitly built its federal spending projections on the assumption that SSI would remain unavailable to Puerto Rico residents (\href{https://www.cbo.gov/system/files/2023-05/51313-2023-05-ssi.pdf}{CBO SSI Baseline, May 2023}). Although territorial delegates repeatedly reintroduced the \emph{SSI Equality Act}---including a renewed push in October 2025---Congress through 2026 has not enacted legislation extending SSI to Puerto Rico, Guam, the U.S. Virgin Islands, or American Samoa (\href{https://puertoricoreport.com/territorial-delegates-reintroduce-ssi-equality-act/}{Puerto Rico Report, Oct. 21, 2025}).

Doctrinally, however, the separate opinions in \emph{Vaello Madero} proved almost as consequential as the 8--1 holding. First, Justice Gorsuch's concurrence---joined in spirit by Justice Sotomayor's Footnote 4---forged a rare cross-ideological call to overrule the \emph{Insular Cases}. Yet when an immediate vehicle arrived months later in \emph{Fitisemanu v. United States} (challenging the Tenth Circuit's reliance on the \emph{Insular Cases} to deny birthright citizenship to persons born in American Samoa), the Supreme Court denied certiorari in October 2022 without recorded dissent (\href{https://lawreview.syr.edu/supreme-court-declines-to-hear-constitutional-challenge-to-citizenship-law/}{Syracuse Law Review, Oct. 27, 2022}). Through 2026, the \emph{Insular Cases} remain formally unoverruled, even as lower courts and litigants increasingly avoid relying on them. Second, Justice Thomas's concurrence reopened foundational debates over \href{https://supreme.justia.com/cases/federal/us/347/497/}{\emph{Bolling v. Sharpe}} and the original meaning of the Fourteenth Amendment's Citizenship Clause---a thread that continues to surface in constitutional scholarship and originalist challenges to federal classification doctrines.

\section{Pedagogical Critiques and Exam Hypotheticals}
\textbf{Pedagogical Critiques.} First, \emph{Vaello Madero} illustrates how ``toothless'' rational-basis review operates at the level of aggregate jurisdiction rather than individual fit: even though no SSI-eligible resident of New York pays federal income tax either, the Court permitted Congress to use Puerto Rico's \emph{territorial} tax exemption as a proxy for benefits eligibility. Second, the majority's Footnote 1---disclaiming any implication that Congress could exclude residents of an individual State---is undertheorized: if Fifth Amendment rational-basis review permits tying benefits to aggregate fiscal contributions, why could Congress not exclude low-revenue States, absent an unstated structural principle of equal state footing or political-process representation (cf. \emph{Carolene Products} Footnote 4)? Third, Justice Gorsuch's concurrence exposes an ironic tension: overruling the \emph{Insular Cases} would not have changed the outcome in \emph{Vaello Madero} itself, because the majority assumed Fifth Amendment equal protection applied in full and still upheld the statute under the Territory Clause.

\textbf{Exam Hypothetical and ``Opposite Side'' Contrast.}
\begin{itemize}
  \item \textbf{Scenario A (Territorial Asymmetry Upheld Under \emph{Vaello Madero}):} Congress enacts a new federal ``National Heating and Cooling Credit'' funded from general Treasury revenues and administered directly by the IRS, available to low-income households in the 50 States and D.C., but excluding residents of Guam and Puerto Rico while maintaining separate block grants for those territories. A disabled U.S. citizen living in Guam sues under the Fifth Amendment. \emph{Analysis:} Congress wins under \emph{Vaello Madero}, \href{https://supreme.justia.com/cases/federal/us/435/1/}{\emph{Torres}}, and \href{https://supreme.justia.com/cases/federal/us/446/651/}{\emph{Rosario}}. Under the Territory Clause and deferential rational-basis review, Guam's distinct federal tax structure supplies a rational basis for territorial exclusion without any dollar-for-dollar accounting.
  \item \textbf{Scenario B (``Opposite Side'' Contrast---Individual State Exclusion or Suspect Classification):} Now suppose Congress either (1) excludes residents of West Virginia from the same federal credit because West Virginia residents pay less per capita in federal income taxes than residents of New York, or (2) excludes Puerto Rico residents from SSI only if they are of Hispanic ancestry while granting SSI to non-Hispanic residents of Puerto Rico. \emph{Analysis:} Both measures fail. In (1), Congress cannot rely on the Territory Clause, and \emph{Vaello Madero} Footnote 1 expressly refuses to extend territorial fiscal-balance deference to exclusions of individual States---which would also collide with structural principles of equal state sovereignty and interstate travel. In (2), the classification ceases to be a geographic/territorial line and becomes an explicit racial or national-origin classification triggering strict scrutiny under Fifth Amendment equal protection (\href{https://supreme.justia.com/cases/federal/us/515/200/}{\emph{Adarand}}; or, on Justice Thomas's view, violating the Citizenship Clause), which the federal government cannot survive.
\end{itemize}

\end{document}
